\documentclass[11pt]{article}
\usepackage[margin=1in]{geometry}
\usepackage{amsmath,amssymb,amsthm,mathtools}
\usepackage[hidelinks]{hyperref}
\setlength{\emergencystretch}{3em}
\allowdisplaybreaks[2]
\newtheorem{proposition}{Proposition}[section]
\newtheorem{lemma}[proposition]{Lemma}
\newtheorem{corollary}[proposition]{Corollary}
\theoremstyle{remark}\newtheorem{remark}[proposition]{Remark}
\newcommand{\E}{\mathbb E}
\newcommand{\dd}{\,d}
\newcommand{\one}{\mathbf1}
\newcommand{\cR}{\mathcal R}
\title{Distance control on the c-form chart and signed radial gain\\
S2a: a partial stopped implication and the extension obstruction}
\author{Analytical increment for review}
\date{Version 1; October 21, 2026}
\begin{document}
\maketitle
\begin{abstract}
A Gaussian change of the input coordinate proves the requested
distance rate on the fixed c-form chart, without an integral
assumption on the reference offset or the common forcing.
For every original-flow receiver we also prove a signed radial
upper bound with integrated error $4B+o(1)$. This includes
the lower-side weak gate and gives the capped mass gain as soon
as the distance envelope is available up to physical exit.
The full requested implication is not closed: the fixed-chart
hypothesis does not provide the off-chart pair comparison used
between scaled-chart exit and fixed physical exit. We display
the unpaid inequality and the precise failure of substituting
the existing full-gate estimates for it. The secondary energy
and complement mass conclusions are consequently not asserted.
\end{abstract}

\section{Scope, normalization, and stopped assumptions}
All new labels have prefix \texttt{an03decgv1:}.
Use the exact original balanced flow and label measure of CLOCK,
TAILREP and LCRC:
\[
 P=\tfrac12P_1+\tfrac12P_2,\quad
 P_1=N(e_1,q^2I_2),\quad P_2=N(\rho e_2,q^2I_2),\quad
 \rho\in[13/20,7/10],\quad \lambda=\rho^2,\quad Q=\lambda+q^2,
\]
\[
 U'=R_\theta^TW,\quad W'=R_\theta U,\quad
 R_\theta=\E_P[(X-f)X^T\one_{\{a_\theta\cdot X>0\}}],
 \quad U=\sqrt m\,a_\theta,\quad W=\sqrt m\,a_\psi.
\]
Write $e_j=X_j-f_j$, $c=Q^{-1}\E_{P_2}[X_2f_2]$,
$I(t)=Q\int_{t_0}^t(1-c)$, $\beta=\lambda/Q$ and
$\ell=\log(1/q)$. Here $t_0=t_h$ and
$t_*-t_0\le L_T\ell$, as in the prescribed LCRC interval.
The original reference is retained throughout. Its scaled
coordinates are $(x_c,y_c)=(\theta_c/q,\psi_c/q)$.
Its radial variable $m_c$ solves the same radial equation; it
need not represent an atom of positive population mass.

Assume the question's stopped guards, in particular
\begin{equation}
 0\le c\le c_b<1,\quad
 r(t):=\sup_\theta|R_{11}(t,\theta)|,\quad \int_{t_0}^{t_*}r<B,
 \quad \|e_2\|_{P_1}+\|e_1\|_{P_2}\le M_{\rm res}q.
 \label{an03decgv1:guards}
\end{equation}
The notation $M_{\rm res}$ distinguishes the residual guard from
strong mass. Loss monotonicity gives
$\|e_1\|_{P_1}+\|e_2\|_{P_2}\le C_{\mathcal L}$, with a
fixed initial-loss constant. The other stopped guards
$Y<K_Y$ and $E_{G_+}<E_g$ are retained but are not needed
in the two new estimates below.

On $[-L_c,L_c]^2$ assume exactly CLOCK's c-form contract:
\begin{align}
 2f_x&=-(1-M_s)x+\rho\phi(\rho x)-\rho\mathcal S(x)
       +\lambda[V+(1-c)y]\Phi(\rho x)+2\cR_x,\nonumber\\
 2f_y&=-(1-M_s)y-Y_s-K_w+\rho(1-c)\mathcal K(\rho x)
       +2\cR_y,\label{an03decgv1:cform}\\
 \mathcal S(x)&=\int F(\rho x,\rho r_0)\dd\mu_s(r_0,y_0),\quad
 F(z,w)=\phi(\min(z,w))+w\Phi(\min(z,w)),\nonumber\\
 \mathcal K(z)&=\phi(z)+z\Phi(z),\quad
 \mu_s\ge0,\quad \mu_s(\mathbb R^2)=M_s\le1.\nonumber
\end{align}
The mass restriction is part of CLOCK's stated c-form contract,
not a conclusion about the population in this increment.
The donor measure and $V,Y_s,K_w,c,M_s$ are held fixed when
differentiating with respect to a receiver. Assume
\[
 |\cR|_1\le\epsilon_F,\quad
 \|D_{x,y}\cR\|_{1\to1}\le\epsilon_F,\qquad
 \int_{t_0}^{t_*}\epsilon_F\le F_b.
\]
Changing equivalent value-norm conventions only enlarges a fixed
constant. $L_c\ge79/51$ is fixed before the residual guards.
CLOCK's offset guard, with its allowed $e_*\le1$, and
$a_\rho<28/51$ imply
\begin{equation}
 |x_c|,|y_c|\le R_c:=79/51\le L_c.
 \label{an03decgv1:reference}
\end{equation}
We do not apply CLOCK's passage theorem or improve this guard.

Use the fixed core/complement ancestry of TAILREP, and SEP's
$\Sigma=\{c_0<Kq^\gamma\}$, with $\mathcal D=\emptyset$.
For a receiver define, using continuous angle lifts,
\[
 D_\alpha(t)=\frac{|\theta_\alpha-\theta_c|
                         +|\psi_\alpha-\psi_c|}{q},\quad
 \xi_\alpha=D_\alpha(t_0),\quad D_{0,\alpha}=q\xi_\alpha.
\]
Thus $\xi$ is scaled displacement, whereas $D_0$ in the
gain-integration formula below is physical displacement.
No label is removed when it leaves a chart.

\section{A distance envelope on the fixed c-form chart}
\begin{proposition}[A coordinate change removes the common forcing]
\label{an03decgv1:local}
Let a receiver and the reference stay in $[-L_c,L_c]^2$
on $[s,t]$. Under \eqref{an03decgv1:cform},
\begin{equation}
 D_\alpha(t)\le C_{D,\mathrm{loc}}D_\alpha(s)
        \exp\!\left(\frac\beta2[I(t)-I(s)]\right),\qquad
 C_{D,\mathrm{loc}}=p_*^{-1}e^{C_TF_b}.
 \label{an03decgv1:localdistance}
\end{equation}
Here $p_*:=\Phi(-7L_c/10)>0$ and one may take
$C_T=2+p_*^{-1}+h_*$, where
$h_*=(7/10)\phi(0)/p_*$.
These constants are independent of $B,M_{\rm res},K_Y,E_g$,
and of the values of $V,Y_s,K_w$.
\end{proposition}
\begin{proof}
Put $p(x)=\Phi(\rho x)$ and change only the input coordinate:
\[
 u=\mathcal T_\rho(x):=\int_0^x\frac{\dd v}{p(v)},
 \qquad \widetilde f_u=f_x/p(x),\qquad \widetilde f_y=f_y.
\]
For the leading field, the $V$ term is the additive constant
$\lambda V/2$. The off-diagonal derivatives are exactly
\begin{equation}
 \partial_y\widetilde f_u=\lambda(1-c)/2,\qquad
 \partial_u\widetilde f_y=\lambda(1-c)p(x)^2/2,
 \qquad \partial_y\widetilde f_y=-(1-M_s)/2\le0.
 \label{an03decgv1:offdiagonals}
\end{equation}
To check the remaining diagonal, decompose the scalar part as
\begin{equation}
 \frac{-(1-M_s)x+\rho\phi(\rho x)-\rho\mathcal S(x)}{2p(x)}
 =(1-M_s)\frac{-x+\rho\phi(\rho x)}{2p(x)}
 +\int\frac{\rho[\phi(\rho x)-F(\rho x,\rho r_0)]}{2p(x)}
       \dd\mu_s.
 \label{an03decgv1:decomposition}
\end{equation}
For $z=\rho x$ and $h=\phi(z)/\Phi(z)$, differentiation
with respect to $u$ of the free term gives
\[
 -\tfrac12\{1-zh+\lambda\phi(z)(z+h)\}<0.
\]
Indeed $z+h=\mathcal K(z)/\Phi(z)>0$. For $z<0$, $1-zh>0$;
for $z\ge0$, $\Phi(z)-z\phi(z)>0$, since its value at zero
is $1/2$ and its derivative is $z^2\phi(z)\ge0$.

For one donor, the last integrand in
\eqref{an03decgv1:decomposition} is constant $-\lambda r_0/2$
when $x\le r_0$. When $x>r_0$, its $u$ derivative is
\[
 \frac{\lambda\phi(\rho x)}{2p(x)}
       [\mathcal K(\rho r_0)-\mathcal K(\rho x)]\le0,
\]
because $\mathcal K'=\Phi>0$. The one-sided derivatives agree
at equality. Positivity of $\mu_s$ and $M_s\le1$ therefore
give $\partial_u\widetilde f_u\le0$.
The leading transformed Jacobian is Metzler and each column
sum is at most $\lambda(1-c)/2$. Its induced $\ell^1$
logarithmic norm has that same bound. This sign assertion is
for the leading field only; no finite-noise cooperativity is assumed.

The transformed remainder is $(\cR_x/p,\cR_y)$. Its derivatives
are
\[
 \begin{pmatrix}
  \partial_x\cR_x-(p'/p)\cR_x&\partial_y\cR_x/p\\
  p\partial_x\cR_y&\partial_y\cR_y
 \end{pmatrix}.
\]
Since $p_*\le p\le1$ and $0\le p'/p\le h_*$ on the chart,
its induced $\ell^1$ norm is at most $C_T\epsilon_F$.
The transformed chart is a rectangle. The segment joining
the two transformed receivers lies in that rectangle, so the
mean-value formula and the logarithmic-norm inequality give
\[
 D^+\bigl(|u-u_c|+|y-y_c|\bigr)
 \le[\lambda(1-c)/2+C_T\epsilon_F]
                 \bigl(|u-u_c|+|y-y_c|\bigr).
\]
Finally $1\le\mathcal T_\rho'\le p_*^{-1}$ and
$\lambda\int(1-c)=\beta\Delta I$. Gronwall proves
\eqref{an03decgv1:localdistance}, including coincident receivers.
All estimates are instantaneous in the donor measure; no donor
weight or complement moment is differentiated.
\end{proof}

\section{A global signed radial estimate}
\begin{lemma}[Projected gates with a radial weight]
\label{an03decgv1:projected}
Under \eqref{an03decgv1:guards}, for every receiver let
$\sigma=\sin\theta$ and $d_2=Q(1-c)/2=I'/2$.
There are constants $C,k>0$, depending only on the ratio box
and initial loss, such that
\begin{align}
 |R_{22}-d_2|&\le CM_{\rm res}q^2+Ce^{-k\sigma^2/q^2}
       &&(\sigma\ge0),\label{an03decgv1:positivegate}\\
 |R_{22}|&\le CM_{\rm res}q^2+Ce^{-k\sigma^2/q^2}
       &&(\sigma<0).\label{an03decgv1:negativegate}
\end{align}
Consequently, for all output directions, without a quadrant
assumption,
\begin{equation}
 2R_{22}\sin\theta\sin\psi
 \le I'(\sin\theta\sin\psi)_++Cq+CM_{\rm res}q^2.
 \label{an03decgv1:signedweak}
\end{equation}
\end{lemma}
\begin{proof}
The $P_1$ part of $R_{22}$ is bounded by
$\frac12\|e_2\|_{P_1}\|X_2\|_{P_1}\le M_{\rm res}q^2/2$.
On $P_2$, $a_\theta\cdot X$ has mean $\rho\sigma$ and
variance $q^2$. The wrong-side probability is at most
$\exp(-\rho^2\sigma^2/(2q^2))$. For that side $A$,
\[
 |\E_{P_2}[e_2X_2\one_A]|
 \le\|e_2\|_{P_2}(\E_{P_2}X_2^4)^{1/4}
                         \mathbb P_{P_2}(A)^{1/4}
 \le C e^{-k\sigma^2/q^2},\qquad k=(13/20)^2/8.
\]
For $\sigma\ge0$, subtract the lost gate from the exact
global moment $\E_{P_2}[e_2X_2]=Q(1-c)$.
For $\sigma<0$, bound the selected gate itself. This proves
\eqref{an03decgv1:positivegate}--\eqref{an03decgv1:negativegate}.
Multiplication by $|\sin\theta\sin\psi|\le|\sigma|$ costs
at most $Cq$, because
$\sup_{v\ge0}v e^{-kv^2/q^2}=q/\sqrt{2ke}$.
Since $I'\ge0$, replacing $\sin\theta\sin\psi$ by its
positive part is an upper bound. This proves
\eqref{an03decgv1:signedweak} on both sides of the weak gate.
\end{proof}

\begin{proposition}[Radial bound with the actual displacement]
\label{an03decgv1:radial}
Uniformly for every label and every endpoint of the stopped interval,
\begin{equation}
 \left(\log\frac{m_\alpha}{m_c}\right)'
 \le I'\min\{1,q^2D_\alpha(t)^2\}+b_q(t),\qquad
 b_q(t)=4r(t)+C(1+M_{\rm res})q,
 \label{an03decgv1:actualgain}
\end{equation}
where $C$ depends only on $C_{\mathcal L}$, the ratio box and
$R_c$. In particular
\begin{equation}
 \int_{t_0}^t b_q\le4B+C(1+M_{\rm res})qL_T\ell=4B+o(1),
 \qquad \frac{m_c(t)}{m_c(t_0)}\le e^{2B+o(1)}.
 \label{an03decgv1:radialbudget}
\end{equation}
These bounds do not require a c-form representation away from
the reference, a lower-side projected-margin evolution, or a
sign for $R_{11}$.
\end{proposition}
\begin{proof}
For any selector, the residual norms and Gaussian second moments
give
\[
 |R_{12}|+|R_{21}|\le C(1+M_{\rm res})q,\qquad
 |R_{22}|\le C+CM_{\rm res}q^2.
\]
The exact balanced radial identity is
$(\log m)'=2a_\psi^TR_\theta a_\theta$.
The two strong-diagonal terms in the difference between a label
and the reference are bounded in absolute value by $4r$;
the cross terms by $C(1+M_{\rm res})q$.
By \eqref{an03decgv1:reference}, the reference weak radial
term is bounded in absolute value by $Cq^2+CM_{\rm res}q^4$.
Apply Lemma \ref{an03decgv1:projected} to the other weak term.

For completeness, TAILREP's scalar $r_1$ is arbitrary in the
exact splitting $R_\theta=\operatorname{diag}(r_1,d_2)
+\mathcal B_\theta$. Taking $r_1=R_{11}(\theta_c)$, its
$11$ splitting defect must be recombined with the displayed
$2r_1$ term. Their sum is exactly
\[
 2R_{11}(\theta)a_1b_1
      -2R_{11}(\theta_c)a_1^cb_1^c,
\]
which is bounded by $4r$. Bounding both split contributions
separately is unnecessary; omitting the $11$ defect would be
incorrect.

The elementary Lipschitz bound for sine gives
\[
 (\sin\theta\sin\psi)_+
 \le\tfrac12(\sin^2\theta+\sin^2\psi)
 \le q^2D_\alpha^2+2R_c^2q^2.
\]
It is also at most one. Use
$\min(1,A+b)\le\min(1,A)+b$ and $0\le I'\le Q$;
absorb the resulting $O(q^2)$ term into $b_q$.
This proves \eqref{an03decgv1:actualgain} even after angular
transit. Its integrated bound follows from the duration guard.
For the reference alone the cross products are $O(q)$ and
the weak product is $O(q^2)$, so
$(\log m_c)'\le2r+C(1+M_{\rm res})q^2$.
This gives the last bound in \eqref{an03decgv1:radialbudget}.
\end{proof}

\begin{corollary}[The remaining premise for capped gain is distance]
\label{an03decgv1:conditionalgain}
For a specified label, suppose in addition that
\begin{equation}
 D_\alpha(t)\le C_D\xi_\alpha e^{\beta I(t)/2}
 \quad\text{until its first }qD_\alpha=\delta_0,
 \label{an03decgv1:missingdistance}
\end{equation}
where $0<\delta_0<1$ is fixed. This premise has only been proved
here while the receiver stays in the fixed scaled chart.
Set $A_0=C_D^2\max(1,\delta_0^{-2})$, increasing $C_D$ to at
least one. Then at every endpoint, including after physical exit,
\begin{align}
 \left(\log\frac{m_\alpha}{m_c}\right)'
 &\le I'\min\{1,A_0D_{0,\alpha}^2e^{\beta I}\}+b_q,
 \label{an03decgv1:capped}\\
 \frac{m_\alpha(t)}{m_\alpha(t_0)}
 &\le e^{6B+1/\beta+o(1)}
       \max\{1,A_0q^2\xi_\alpha^2e^{\beta I(t)}\}^{1/\beta}.
 \label{an03decgv1:massgain}
\end{align}
The exponent is $1/\beta=1+q^2/\lambda\le1+1/100$ when
$q^2\le\lambda/100$.
\end{corollary}
\begin{proof}
Before exit use \eqref{an03decgv1:missingdistance} in
\eqref{an03decgv1:actualgain}. At exit continuity gives
$C_D^2q^2\xi_\alpha^2e^{\beta I}\ge\delta_0^2$.
Monotonicity of $I$ makes the minimum in
\eqref{an03decgv1:capped} equal to one at all subsequent times,
even if the label reenters. A label initially past the physical
cutoff is handled in the same way at entry.
Apply TAILREP's already proved gain-integration lemma at each
endpoint, then multiply by the reference gain in
\eqref{an03decgv1:radialbudget}. This proves
\eqref{an03decgv1:massgain}. In this radial estimate the geometric
cost is $\delta_0^{-2}$ in $A_0$, not an extra gain exponent.
Any loss in an eventual proof of the missing distance premise
would still have to be retained; it has not been set to zero here.
\end{proof}

\section{Exact stopping point: fixed scaled chart versus physical exit}
\label{an03decgv1:obstructionsection}
For an initially charted label let $\tau_L$ be its first exit
from $[-L_c,L_c]^2$ and $\tau_\delta$ its first physical exit.
At the former time
\[
 qD_\alpha(\tau_L)\le2q(L_c+R_c)=O(q),
\]
which is smaller than $\delta_0$ for small $q$. Proposition
\ref{an03decgv1:local} therefore does not cover
$[\tau_L,\tau_\delta]$. Some of SEP's labels can also start
outside the specified square: SEP supplies a fixed cap
$Z_\Sigma(K)$, not the assertion that this cap fits the prescribed
$L_c$.

The exact comparison still needed for a chart-exiting label is
\begin{equation}
 \boxed{\quad
 \sup_{\tau_L\le t\le\tau_\delta\wedge t_*}
 \frac{D_\alpha(t)}{D_\alpha(\tau_L)}
 \exp\!\left[-\frac\beta2(I(t)-I(\tau_L))\right]
 \le C_{\rm ext}e^{CB}.
 \quad}
 \label{an03decgv1:unpaid}
\end{equation}
Here $C_{\rm ext}$ must be fixed independently of
$M_{\rm res},K_Y$ and $q$; fixed chart and remainder-budget
dependence is allowed. An initially exterior label needs the
corresponding comparison starting at $t_0$. Neither comparison
is added as a hypothesis of a claimed full S2a theorem.

\paragraph{What the off-chart calculation supplies, and where it fails.}
Adding weak output to the full-gate calculation does give
\eqref{an03decgv1:positivegate}. It gives a matrix \emph{value}
estimate, not the required same-field distance comparison.
If a positive receiver has only a fixed projected margin
$\sin\theta/q\ge L>0$, using its exponential tail by absolute
value in a Lipschitz comparison costs at best
\begin{equation}
 \int_{\tau_L}^{t_*}Ce^{-kL^2}\dd t
 \le CL_Te^{-kL^2}\ell.
 \label{an03decgv1:fixedmargin}
\end{equation}
For fixed $L$ this bound is not uniform as $q\downarrow0$.
It cannot be substituted for a fixed error budget in
\eqref{an03decgv1:unpaid}. A growing projected margin, a signed
cancellation, or another off-chart pair estimate must be proved.
The factor $|\sin\theta|$ that paid this tail in the radial
lemma does not automatically occur in an angular derivative.
On the lower side even the leading full-gate value is different:
\eqref{an03decgv1:negativegate} replaces
\eqref{an03decgv1:positivegate}.

More explicitly, differentiation of the selected residual matrix
at a fixed population state gives the boundary moment
\begin{equation}
 \partial_\theta R_\theta
 =\E_P\!\left[(X-f)X^T(a_\theta^\perp\!\cdot X)
                         \delta(a_\theta\cdot X)\right].
 \label{an03decgv1:boundary}
\end{equation}
The formula may equivalently be written as a Gaussian integral
over the projection hyperplane. An $L^2$ residual bound is not
by itself a trace bound on that hyperplane. The c-form derivative
guard estimates its contribution on the fixed square only.
This is the precise derivative gap in trying to use a full-gate
value estimate as an off-chart logarithmic-norm estimate.

Q2EC's full-gate and transport lemmas do not fill that gap after
$t_h$. Their proof uses, on $[t_m,t_h]$, strong donors in
$[-R,R]^2$, $M\le1/2$, $M'/M\ge1/3$,
$\mu\le Cq^2\ell^2M$, and integrated donor-spread and outside
budgets. Those time-restricted controls cannot be extended past
$t_h$ by changing the endpoint in the lemma. In particular its
integrated $O(M/x)$ transport error uses this structured mass
and angular evolution. Substituting the residual-guard constant
$M_{\rm res}$ for that mass proves neither the requisite
integrability nor guard-independent constants.

Finally, directly extending the coordinate-change proof would
require \eqref{an03decgv1:cform}, including its derivative
remainder budget, on the receiver region out to distance
$\delta_0/q$. That is an additional field hypothesis, not the
given fixed-chart hypothesis. The result is stopped at
\eqref{an03decgv1:unpaid}; no such extension is assumed.
This is a failure to close the stated implication from the
available estimates, not a counterexample to the initialized flow.

\section{Consequences not claimed, outside ancestry, and constants}
\paragraph{Secondary and strong-complement contracts.}
SEP supplies LCRC's entry-profile contract, and MOMC supplies
the strong entry profile. Neither is reproved. The new radial
bound closes the signed radial-defect part of S2 and reduces its
gain contract to \eqref{an03decgv1:missingdistance}.
It does not establish that missing contract for every label.
Consequently BULK Proposition 5.2 and its secondary corollary,
and TAILREP's tail-moment lemma, cannot yet be applied to claim
\[
 E_\Sigma(t)\le C_\Sigma(q^5+q^{\Gamma_\nu}H(t)),\qquad
 \int_{t_0}^{t_*}\mu_{\rm comp}/q\le Cq^{\Delta_\nu}\ell.
\]
Their already proved conditional margins remain
$\Gamma_\nu\ge289169/924000>3/10$ and
$\Delta_\nu\ge17141311/140024000>3/25$ for $\nu\le1/100$.
These rational margins are not new original-flow conclusions.
No extra displacement-growth charge is inserted into the
mass-charged tail calculation.

\paragraph{Q3 leakage is a separate unpaid outside contribution.}
The initial Q3 interval $(\pi,3\pi/2)$ is neither MOMC's
right-half strong ancestry nor SEP's Q2 group. It cannot be
included in either entry tail by a change of name. The available
coarse MOMC left-ancestry estimate gives only
$\mu_{Q3}(t_0)\le Cq^2\ell^2$. The same proof as the radial
lemma, without subtracting the reference, gives
\[
 \mu_{Q3}(t)\le\mu_{Q3}(t_0)e^{2B+I(t)+o(1)}.
\]
Even if one uses the stopped clock estimate from BULK and the
seed $H(t_0)\asymp A_q^{-\lambda}q^k$, $k=10(1-\lambda)$,
this rough estimate gives only
\[
 \int_{t_0}^{t_*}\mu_{Q3}/q
 \le C e^{2B}E_g A_q^{\lambda}q^{1-k}\ell^3.
\]
Since $k\ge51/10$, this upper bound does not tend to zero.
It is not a lower bound on the actual leakage. For this ancestry,
including the neighborhoods of both $\pi$ and $3\pi/2$, one
still needs a proved stopped integrated-mass bound
$\int\mu_{Q3}/q=o(1)$, or a sharper entry/gain argument implying
it. Neither entry tail used above supplies that argument.

\paragraph{Constant audit.}
The fixed constants are the ratio box, initial-loss bound,
$L_T,L_c,R_c$, the positive Gaussian minimum $p_*$,
$h_*,C_T,k$, and the remainder budget $F_b$.
The local distance constant is exactly
$p_*^{-1}e^{C_TF_b}$. The radial error integral is
$4B+C(1+M_{\rm res})qL_T\ell$; the reference gain is
$e^{2B+C(1+M_{\rm res})q^2L_T\ell}$.
Conditional on the unpaid envelope, the mass-gain constant is
$e^{6B+1/\beta+o(1)}$ and
$A_0=C_D^2\max(1,\delta_0^{-2})$.
Thus $B$ appears only through a fixed exponential and the
displayed additive budget; $M_{\rm res}$ only multiplies
vanishing scales. There is no $K_Y$ dependence in the proved
lemmas. $E_g,A_q,K,Z_\Sigma$ occur only in the cited entry/clock
audit, not in the local distance or radial constants.
The missing $C_{\rm ext}$ has not been declared finite.

\begin{thebibliography}{9}\small
\bibitem{ROAD} Unified analytic theory roadmap, revision 1.9,
S0 assembly, S2 October 21 update, S3 status and ledger 16.
\bibitem{TAIL} TAILREP, \texttt{AN03\_mass\_charged\_tail\_and\_selection\_defects\_v1}:
gain integration, tail moment, rational margins, and final exact radial identity.
\bibitem{BULK} BULK, \texttt{AN03\_bulk\_clock\_and\_secondary\_Q2\_energy\_v1}:
stopped clock, Lemma 5.1, Proposition 5.2 and secondary corollary.
\bibitem{CLOCK} CLOCK, \texttt{AN03\_exact\_amplitude\_clock\_and\_c\_passage\_v1}:
exact response clock and c-form contract.
\bibitem{LCRC} LCRC, \texttt{AN03\_ledger\_cap\_residual\_closure\_v1},
Section 1 contracts and stopped clock lemma.
\bibitem{CFR} CFR, \texttt{AN03\_c\_form\_field\_reduction\_v1}:
reviewed local source; selected moments and outstanding S2 dependencies.
\bibitem{ENTRY} Q2EC, SEP and MOMC: original entry contracts,
recruited profile, and moving-overlap strong passage. Exact paths
and source hashes are recorded in the accompanying review note.
\bibitem{PROFILE} PROFILE, \texttt{AN02\_strong\_entry\_profile\_and\_tail\_budget\_v1},
Proposition 4.1: outside output costs $\mu/q$, including receiver derivative.
\end{thebibliography}
\end{document}
